\section{基于结构方程模型的噪声影响传导机制}

有序Logistic回归揭示了各因子对声环境满意度的直接效应，但未刻画噪声源因子与身心、情绪因子之间的传导关系。为理解噪声影响的形成路径，本节构建结构方程模型进行进一步检验。

\subsection{模型设定}

结构方程模型由测量模型与结构模型两部分组成。测量模型将四个潜变量分别对应探索性因子分析中的题项：$F_1$ 对应入睡、注意力、精神、听觉四项，$F_2$ 对应睡眠中断、人际关系、身体、情绪四项，$F_3$ 对应交通、施工、公共活动三项，$F_4$ 对应生活、商业、动物三项。结构模型设定噪声源因子 $F_3$、$F_4$ 作用于身心因子 $F_1$ 与情绪因子 $F_2$，身心与情绪因子进一步作用于声环境满意度，同时允许 $F_1$ 与 $F_2$ 相关。

\subsection{模型拟合与路径分析}

模型拟合指标为 $\chi^2=261.205$（df=81）、CFI=0.916、TLI=0.891、RMSEA=0.069，整体拟合可以接受。结构路径的标准化系数如表\ref{tab:sem}所示。

\begin{table}[H]
\centering
\caption{结构方程模型的标准化路径系数}
\label{tab:sem}
\rowcolors{2}{tablight}{white}
\begin{tabular}{M{7.5cm}M{4cm}M{3cm}}
\bluetoprule
\rowcolor{tabhead}
\thead{路径} & \thead{标准化系数} & \thead{显著性} \\
\hline
$F_3$ 环境源 $\to F_1$ 身心 & 0.370 & $<$0.001 \\
$F_4$ 生活源 $\to F_1$ 身心 & 0.072 & 不显著 \\
$F_3$ 环境源 $\to F_2$ 情绪 & 0.216 & $<$0.05 \\
$F_4$ 生活源 $\to F_2$ 情绪 & 0.352 & $<$0.01 \\
$F_1$ 身心 $\to$ 满意度 & 0.003 & 不显著 \\
$F_2$ 情绪 $\to$ 满意度 & $-0.269$ & $<$0.001 \\
$F_3$ 环境源 $\to$ 满意度 & $-0.214$ & $<$0.05 \\
$F_4$ 生活源 $\to$ 满意度 & $-0.277$ & $<$0.05 \\
$F_1 \leftrightarrow F_2$ & 0.188 & $<$0.001 \\
\bluebottomrule
\end{tabular}
\end{table}

路径系数显示，噪声源因子对身心与情绪因子均存在正向传导作用，其中典型环境噪声源对身心因子的效应最强（0.370），生活场景噪声源对情绪因子的效应最强（0.352）。在满意度层面，情绪因子对满意度的直接效应为 $-0.269$（$P<0.001$），是连接噪声源与满意度的核心通道；身心因子对满意度的直接效应不显著，其影响被情绪因子所中介；噪声源因子对满意度亦存在直接的负向效应。

上述结果表明，噪声对满意度的影响包含直接与间接两条路径。间接路径为：噪声源先损害居民身心状态，再经由情绪与社会关系渠道最终作用于满意度。情绪与社会关系因子在其中承担关键的中介作用，这与有序Logistic回归中该因子系数最强的结果相互印证。

